\chapter{Encodings and finite probability}
\label{chap:encodings}

This chapter covers the machine-independent layer that every other chapter
relies on: how objects become binary strings, and how probabilities over random
bit strings are counted. Formalized so far: fixed-length bit strings and their
bridge to lists, the self-delimiting block framing and the pairing function
$\langle x, y \rangle$ used for every two-argument machine input, canonical
binary natural numbers, a rose-tree serialization, a canonical code for
fan-in-two circuits, exact counting over $\{0,1\}^T$, uniform event
probability with the union bound and majority amplification, and a fixed-time
repetition wrapper for probabilistic machines whose acceptance probability is
exactly a majority-event probability. The main open direction is
consolidation: about 24 modules define ad hoc codecs and about 25 define ad hoc
probability notions, which a reusable codec interface
(Definition~\ref{def:encodings-codec}) and Mathlib's \texttt{Finset.expect} and
\texttt{Finset.dens} should absorb.

\section{Bit strings}

\begin{definition}[Bit string]
  \label{def:bitstring}
  \lean{Complexity.BitString}
  \leanok
  A bit string of length $n$ is a function $\{0, \dots, n-1\} \to \{0,1\}$.
  Circuits consume bit strings of a fixed length; machines consume lists.
\end{definition}

\begin{definition}[Bit strings as lists]
  \label{def:encodings-bitstring-list}
  \lean{Complexity.BitString.toList, Complexity.BitString.ofList, Complexity.BitString.toTotal}
  \leanok
  \uses{def:bitstring}
  $\mathrm{toList}(x) = [x_0, \dots, x_{n-1}]$ serializes in increasing index
  order; $\mathrm{ofList}$ reads a list whose length is proved to be $n$; and
  $\mathrm{toTotal}(x) : \mathbb{N} \to \{0,1\}$ extends $x$ by $0$ outside
  its range. The variable-length correspondence
  $\{0,1\}^* \simeq \Sigma_n \{0,1\}^n$ is Mathlib's.
\end{definition}

\begin{lemma}[List bridge]
  \label{lem:encodings-bitstring-roundtrip}
  \lean{Complexity.BitString.length_toList, Complexity.BitString.toList_ofList, Complexity.BitString.ofList_toList, Complexity.BitString.toList_inj, Complexity.BitString.toList_append}
  \leanok
  \uses{def:encodings-bitstring-list}
  $|\mathrm{toList}(x)| = n$; $\mathrm{toList}(\mathrm{ofList}(\ell)) = \ell$
  and $\mathrm{ofList}(\mathrm{toList}(x)) = x$; $\mathrm{toList}$ is
  injective; and $\mathrm{toList}$ sends concatenation of bit strings to list
  append.
\end{lemma}
\begin{proof}
  \leanok
  Mathlib's \texttt{List.ofFn} lemmas.
\end{proof}

\section{Self-delimiting blocks and pairing}

\begin{definition}[Self-delimiting blocks]
  \label{def:encodings-delimit}
  \lean{Complexity.delimit, Complexity.unpair?, Complexity.undelimitBlocks}
  \leanok
  The block of $x = x_1 \cdots x_n$ is
  $\mathrm{delimit}(x) = x_1 x_1\, x_2 x_2 \cdots x_n x_n\, 0 1$: every bit is
  doubled and the separator $01$ ends the block. The parser
  $\mathrm{unpair?}$ reads one block off the front of a string and returns the
  payload and the remaining suffix, or fails on malformed input;
  $\mathrm{undelimitBlocks}$ parses a whole sequence of blocks.
\end{definition}

\begin{lemma}[Block framing round trips]
  \label{lem:encodings-delimit}
  \lean{Complexity.delimit_length, Complexity.unpair?_delimit_append, Complexity.eq_delimit_append_of_unpair?_eq_some, Complexity.undelimitBlocks_flatten_delimit}
  \leanok
  \uses{def:encodings-delimit}
  $|\mathrm{delimit}(x)| = 2|x| + 2$. For all $x, y$,
  $\mathrm{unpair?}(\mathrm{delimit}(x)\, y) = (x, y)$, and conversely if
  $\mathrm{unpair?}(z) = (x, y)$ then $z = \mathrm{delimit}(x)\, y$. For every
  list of strings $b_1, \dots, b_m$,
  $\mathrm{undelimitBlocks}(\mathrm{delimit}(b_1) \cdots \mathrm{delimit}(b_m))
  = [b_1, \dots, b_m]$.
\end{lemma}
\begin{proof}
  \leanok
  Induction on the payload; no doubled pair equals the separator.
\end{proof}

\begin{definition}[Pairing function]
  \label{def:pair}
  \lean{Complexity.pair}
  \leanok
  \uses{def:encodings-delimit}
  $\langle x, y \rangle = \mathrm{delimit}(x)\, y$: the first component is
  framed as a block and the second is appended verbatim.
\end{definition}

\begin{lemma}[Length of a pair]
  \label{lem:encodings-pair-length}
  \lean{Complexity.pair_length}
  \leanok
  \uses{def:pair}
  $|\langle x, y \rangle| = 2|x| + 2 + |y|$.
\end{lemma}
\begin{proof}
  \leanok
  \uses{lem:encodings-delimit}
  Induction on $x$.
\end{proof}

\begin{lemma}[Pairing is injective and decodable]
  \label{lem:encodings-pair-inj}
  \lean{Complexity.pair_inj, Complexity.unpair?_eq_some_iff}
  \leanok
  \uses{def:pair}
  If $\langle x_1, y_1 \rangle = \langle x_2, y_2 \rangle$ then $x_1 = x_2$ and
  $y_1 = y_2$. Moreover $\mathrm{unpair?}(z) = (x, y)$ if and only if
  $z = \langle x, y \rangle$.
\end{lemma}
\begin{proof}
  \leanok
  \uses{lem:encodings-delimit}
  Induction on $x_1$.
\end{proof}

\begin{definition}[Projections]
  \label{def:encodings-pair-proj}
  \lean{Complexity.pairFst, Complexity.pairSnd}
  \leanok
  \uses{def:pair}
  $\pi_1$ reads doubled bits until the separator (returning the bits read so
  far on malformed input), and $\pi_2$ returns the suffix after the first
  block, or the empty string if there is no valid block. Both are total.
\end{definition}

\begin{lemma}[Projections of a pair]
  \label{lem:encodings-pair-proj}
  \lean{Complexity.pairFst_pair, Complexity.pairSnd_pair, Complexity.pairSnd_length_le}
  \leanok
  \uses{def:encodings-pair-proj}
  $\pi_1 \langle x, y \rangle = x$, $\pi_2 \langle x, y \rangle = y$, and
  $|\pi_2(z)| \le |z|$ for every $z$.
\end{lemma}
\begin{proof}
  \leanok
  Induction on $x$.
\end{proof}

\begin{theorem}[Pairing is polynomial-time]
  \label{thm:encodings-pair-fp}
  \lean{Complexity.pairFst_mem_FP, Complexity.pairSnd_mem_FP, Complexity.mem_FP_pair}
  \leanok
  \uses{def:encodings-pair-proj, def:FP}
  $\pi_1, \pi_2 \in \cc{FP}$, and if $a, b \in \cc{FP}$ then
  $z \mapsto \langle a(z), b(z) \rangle \in \cc{FP}$.
\end{theorem}
\begin{proof}
  \leanok
  \uses{thm:cobham}
  The projections are block scanners of Cobham's algebra, and pairing is a
  Cobham-definable combination.
\end{proof}

\section{Other formalized codes}

\begin{definition}[Binary natural numbers]
  \label{def:encodings-binary-nat}
  \lean{Complexity.BinaryNatCode.encode, Complexity.BinaryNatCode.decode?}
  \leanok
  $\mathrm{bin}(v)$ is the minimal binary expansion of $v \in \mathbb{N}$,
  least significant bit first, with $\mathrm{bin}(0)$ empty. The decoder
  accepts a string only if it is the minimal expansion of its value, so
  representations with trailing high zeros are rejected. The code is not
  prefix-free; an enclosing framing layer fixes the field boundary.
\end{definition}

\begin{lemma}[Binary natural numbers form an exact codec]
  \label{lem:encodings-binary-nat}
  \lean{Complexity.BinaryNatCode.decode?_eq_some_iff, Complexity.BinaryNatCode.length_encode, Complexity.BinaryNatCode.encode_injective}
  \leanok
  \uses{def:encodings-binary-nat}
  $\mathrm{decode}(w) = v$ if and only if $w = \mathrm{bin}(v)$;
  $|\mathrm{bin}(v)|$ is the bit length of $v$; and $\mathrm{bin}$ is
  injective.
\end{lemma}
\begin{proof}
  \leanok
  Mathlib's \texttt{Nat.bits} API.
\end{proof}

\begin{definition}[Rose-tree data and its serialization]
  \label{def:encodings-data}
  \lean{Complexity.Data, Complexity.Data.toBits, Complexity.DataEncode, Complexity.DataEncode.bitstringEncode}
  \leanok
  A rose tree is a finite list of rose trees. It is serialized as balanced
  brackets: $0$ opens a node, its children follow in order, and $1$ closes it.
  A type with an injective map into rose trees (instances exist for Booleans,
  natural numbers, lists, options, and products) gets a bit-string encoding by
  composing with this serialization. This layer supports the rose-tree machine
  model.
\end{definition}

\begin{lemma}[Rose-tree serialization round trip]
  \label{lem:encodings-data}
  \lean{Complexity.Data.fromBits_toBits, Complexity.Data.length_toBits, Complexity.DataEncode.bitstringEncode_injective}
  \leanok
  \uses{def:encodings-data}
  A decoder recovers every rose tree from its serialization; the serialization
  of $d$ has length equal to the size of $d$; and the derived bit-string
  encoding of any encodable type is injective.
\end{lemma}
\begin{proof}
  \leanok
  A stack-based parser inverts the bracket serialization.
\end{proof}

\begin{definition}[Circuit code]
  \label{def:encodings-circuit-code}
  \lean{Complexity.CircuitCode.encodeCircuit, Complexity.CircuitCode.evalCode}
  \leanok
  \uses{def:circuit}
  A single-output circuit over the fan-in-two AND/OR basis (with negations on
  gate inputs) is encoded as its gate count in terminated unary followed by
  its internal gates in order and then its output gate, each gate as an
  operation bit, two negation bits, and two absolute input wire indices in
  terminated unary. $\mathrm{evalCode}(N, w, x)$ fails unless $|x| = N$;
  otherwise it decodes $w$ exactly (rejecting trailing bits) and evaluates the
  gates in order, failing on any reference to a later wire.
\end{definition}

\begin{theorem}[Circuit code correctness and size]
  \label{thm:encodings-circuit-code}
  \lean{Complexity.CircuitCode.evalCode_encodeCircuit, Complexity.CircuitCode.encodeCircuit_length_le_size}
  \leanok
  \uses{def:encodings-circuit-code, def:bitstring}
  Let $C$ be a single-output fan-in-two AND/OR circuit on $N \ge 1$ inputs with
  size $s$ (internal and output gates). For every $x \in \{0,1\}^N$,
  $\mathrm{evalCode}(N, \mathrm{code}(C), \mathrm{toList}(x)) = C(x)$, and
  $|\mathrm{code}(C)| \le 1 + s(2(N + s) + 6)$.
\end{theorem}
\begin{proof}
  \leanok
  The decoder is exact on canonical codes, and the array-backed evaluator
  agrees with the typed evaluation gate by gate.
\end{proof}

\section{A reusable codec interface}

The codes above were written one at a time, and many more live inside
individual proofs. The following planned nodes describe the shared interface
that should replace them.

\begin{definition}[Binary codec]
  \label{def:encodings-codec}
  A binary codec for a type $\alpha$ with a size measure $\|\cdot\|$ consists of
  an encoder $e : \alpha \to \{0,1\}^*$, a partial decoder
  $d : \{0,1\}^* \to \alpha \cup \{\bot\}$, the round trip $d(e(a)) = a$,
  canonicity ($d(w) = a$ implies $w = e(a)$), and an explicit size bound
  $|e(a)| \le s(\|a\|)$. The existing rose-tree encodings
  (Definition~\ref{def:encodings-data}) give injectivity only, with no decoder
  at the level of $\alpha$ and no size bound.
\end{definition}

\begin{definition}[Pairing and tagged-sum codecs]
  \label{def:encodings-codec-combinators}
  \uses{def:encodings-codec, def:pair}
  \notready
  From codecs for $\alpha$ and $\beta$, the product codec encodes $(a, b)$ as
  $\langle e_\alpha(a), e_\beta(b) \rangle$, with
  $|e(a, b)| = 2|e_\alpha(a)| + 2 + |e_\beta(b)|$, and the tagged-sum codec
  prefixes one tag bit. Both satisfy the codec laws.
\end{definition}

\begin{lemma}[Codec for bounded indices]
  \label{lem:encodings-fin-codec}
  \uses{def:encodings-codec, def:encodings-binary-nat}
  \notready
  For every $n \ge 1$ there is a codec for $\{0, \dots, n-1\}$ all of whose
  codewords have length $\lceil \log_2 n \rceil$, whose decoder rejects values
  $\ge n$.
\end{lemma}

\begin{definition}[Length-bounded list codec]
  \label{def:encodings-list-codec}
  \uses{def:encodings-codec, def:encodings-delimit}
  \notready
  From a codec for $\alpha$, a codec for lists over $\alpha$ that frames each
  element as a block, with
  $|e([a_1, \dots, a_m])| = \sum_i (2|e(a_i)| + 2)$.
\end{definition}

\begin{definition}[Canonical codes for machine objects]
  \label{def:encodings-machine-codes}
  \uses{def:encodings-codec, def:encodings-codec-combinators, lem:encodings-fin-codec, def:encodings-list-codec, def:models-cfg, def:tm}
  \notready
  Canonical codecs, with explicit length bounds, for finite functions (truth
  tables), machine states, configurations of time-bounded runs, and bounded
  interaction transcripts.
\end{definition}

\section{Finite counting}

\begin{lemma}[Size of the sample space]
  \label{lem:encodings-card}
  \lean{Complexity.card_finArrowBool}
  \leanok
  $|\{0,1\}^T| = 2^T$.
\end{lemma}
\begin{proof}
  \leanok
  Mathlib's cardinality of function types.
\end{proof}

\begin{definition}[Blocks of a random string]
  \label{def:encodings-blocks}
  \lean{Complexity.blockEquiv, Complexity.blocksEquiv, Complexity.blockEventCount}
  \leanok
  \uses{def:bitstring}
  $\{0,1\}^{a+b} \simeq \{0,1\}^a \times \{0,1\}^b$ splits a string into prefix
  and suffix, and $\{0,1\}^{kT} \simeq (\{0,1\}^T)^k$ splits it into $k$ blocks
  in row-major order. For an event $E \subseteq \{0,1\}^T$ and
  $w \in \{0,1\}^{kT}$, $\#_E(w)$ is the number of blocks of $w$ lying in $E$.
\end{definition}

\begin{lemma}[Block independence, counting form]
  \label{lem:encodings-block-count}
  \lean{Complexity.card_filter_block, Complexity.card_blockEventCount_eq}
  \leanok
  \uses{def:encodings-blocks}
  The number of $w \in \{0,1\}^{a+b}$ whose prefix satisfies $P$ and whose
  suffix satisfies $Q$ is $|P| \cdot |Q|$. The number of
  $w \in \{0,1\}^{kT}$ with $\#_E(w) = j$ is
  $\binom{k}{j} |E|^j (2^T - |E|)^{k-j}$.
\end{lemma}
\begin{proof}
  \leanok
  Transport along the block equivalences and count tuples by their number of
  successes.
\end{proof}

\begin{lemma}[Union bound, counting form]
  \label{lem:encodings-union-count}
  \lean{Complexity.card_filter_exists_le, Complexity.exists_good_seed}
  \leanok
  The number of points satisfying some $p_i$, $i \in s$, is at most
  $\sum_{i \in s} |p_i|$. Consequently, if bad sets $B_i \subseteq S$ satisfy
  $\sum_i |B_i| < |S|$, some $s \in S$ lies in no $B_i$.
\end{lemma}
\begin{proof}
  \leanok
  Cardinality of a finite union.
\end{proof}

\begin{definition}[Majority]
  \label{def:encodings-majority}
  \lean{Complexity.popCount, Complexity.majority, Complexity.blockMajority}
  \leanok
  \uses{def:encodings-blocks}
  For $f \in \{0,1\}^k$, $\mathrm{maj}(f) = 1$ if and only if more than $k/2$
  positions of $f$ are $1$ (strict majority). For $w \in \{0,1\}^{kT}$,
  $\mathrm{maj}_E(w) = 1$ if and only if $2\,\#_E(w) > k$.
\end{definition}

\begin{lemma}[Majority counts]
  \label{lem:encodings-majority-count}
  \lean{Complexity.card_majority_eq_true, Complexity.card_majority_eq_false, Complexity.card_blockMajority_eq_false}
  \leanok
  \uses{def:encodings-majority}
  $|\{f : \mathrm{maj}(f) = 1\}| = \sum_{r = \lfloor k/2 \rfloor + 1}^{k} \binom{k}{r}$
  and $|\{f : \mathrm{maj}(f) = 0\}| = \sum_{r=0}^{\lfloor k/2 \rfloor} \binom{k}{r}$.
  For $2r + 1$ blocks,
  $|\{w : \mathrm{maj}_E(w) = 0\}| = \sum_{j=0}^{r} \binom{2r+1}{j} |E|^j (2^T - |E|)^{2r+1-j}$.
\end{lemma}
\begin{proof}
  \leanok
  \uses{lem:encodings-block-count}
  Sum the exact fiber counts.
\end{proof}

\section{Event probability}

\begin{definition}[Event probability]
  \label{def:encodings-eventprob}
  \lean{Complexity.eventProb}
  \leanok
  For a finite event $E \subseteq \{0,1\}^T$,
  $\Pr[E] = |E| / 2^T \in \mathbb{Q}$.
\end{definition}

\begin{lemma}[Basic laws]
  \label{lem:encodings-eventprob-basic}
  \lean{Complexity.eventProb_nonneg, Complexity.eventProb_le_one, Complexity.eventProb_mono, Complexity.eventProb_compl, Complexity.eventProb_empty, Complexity.eventProb_univ}
  \leanok
  \uses{def:encodings-eventprob}
  $0 \le \Pr[E] \le 1$; $E \subseteq F$ implies $\Pr[E] \le \Pr[F]$;
  $\Pr[E^c] = 1 - \Pr[E]$; $\Pr[\emptyset] = 0$ and $\Pr[\{0,1\}^T] = 1$.
\end{lemma}
\begin{proof}
  \leanok
  \uses{lem:encodings-card}
  Cardinality arithmetic.
\end{proof}

\begin{lemma}[Union bound]
  \label{lem:encodings-union-bound}
  \lean{Complexity.eventProb_union_le, Complexity.eventProb_biUnion_le, Complexity.eventProb_biUnion}
  \leanok
  \uses{def:encodings-eventprob}
  $\Pr[E \cup F] \le \Pr[E] + \Pr[F]$, and
  $\Pr[\bigcup_{i \in s} E_i] \le \sum_{i \in s} \Pr[E_i]$ for every finite
  family, with equality when the family is pairwise disjoint.
\end{lemma}
\begin{proof}
  \leanok
  Induction on the index set.
\end{proof}

\begin{lemma}[Markov's inequality]
  \label{lem:encodings-markov}
  \lean{Complexity.eventProb_le_uniformAverage_div}
  \leanok
  \uses{def:encodings-eventprob}
  Let $w : \{0,1\}^T \to \mathbb{Q}_{\ge 0}$ and $\theta > 0$. If
  $w(r) \ge \theta$ for every $r \in E$, then
  $\Pr[E] \le \mathbb{E}_r[w(r)] / \theta$, where $\mathbb{E}_r$ is the uniform
  average.
\end{lemma}
\begin{proof}
  \leanok
  Compare sums.
\end{proof}

\begin{lemma}[Block independence]
  \label{lem:encodings-eventprob-block}
  \lean{Complexity.eventProb_block}
  \leanok
  \uses{def:encodings-eventprob, def:encodings-blocks}
  $\Pr_{w \in \{0,1\}^{a+b}}[P(\text{prefix}) \wedge Q(\text{suffix})] =
  \Pr[P] \cdot \Pr[Q]$.
\end{lemma}
\begin{proof}
  \leanok
  \uses{lem:encodings-block-count}
  Divide the counting form by $2^{a+b}$.
\end{proof}

\begin{lemma}[Exact majority tail]
  \label{lem:encodings-majority-tail}
  \lean{Complexity.eventProb_blockMajority_eq_false}
  \leanok
  \uses{def:encodings-eventprob, def:encodings-majority}
  With $p = \Pr[E]$, the probability over $w \in \{0,1\}^{(2r+1)T}$ that
  $\mathrm{maj}_E(w) = 0$ is
  $\sum_{j=0}^{r} \binom{2r+1}{j} p^j (1-p)^{2r+1-j}$.
\end{lemma}
\begin{proof}
  \leanok
  \uses{lem:encodings-majority-count}
  Normalize the exact count.
\end{proof}

\begin{theorem}[Majority amplification]
  \label{thm:encodings-amplification}
  \lean{Complexity.eventProb_blockMajority_false_le_two_pow, Complexity.eventProb_blockMajority_true_ge_one_sub_two_pow, Complexity.eventProb_blockMajority_true_le_two_pow}
  \leanok
  \uses{def:encodings-eventprob, def:encodings-majority}
  Let $E \subseteq \{0,1\}^T$ and $K = 12k + 1$ blocks. If $\Pr[E] \ge 2/3$, then
  $\Pr[\mathrm{maj}_E = 0] \le 2^{-k}$ and $\Pr[\mathrm{maj}_E = 1] \ge 1 - 2^{-k}$.
  If $\Pr[E] \le 1/3$, then $\Pr[\mathrm{maj}_E = 1] \le 2^{-k}$.
\end{theorem}
\begin{proof}
  \leanok
  \uses{lem:encodings-majority-tail, lem:encodings-eventprob-basic}
  Bound the lower tail with $p \ge 2/3$ by $\frac13 (8/9)^r$ using that the
  lower-half binomial coefficients sum to $4^r$ (an elementary bound, not a
  Chernoff bound), take $r = 6k$, and pass to complements for the other two
  statements.
\end{proof}

\begin{lemma}[Good seeds]
  \label{lem:encodings-good-seed}
  \lean{Complexity.exists_good_seed_of_sum_eventProb_lt_one, Complexity.exists_good_seed_of_eventProb_le_two_pow_succ}
  \leanok
  \uses{def:encodings-eventprob}
  If $\sum_{i \in I} \Pr[B_i] < 1$ for bad events $B_i \subseteq \{0,1\}^S$,
  some seed lies in no $B_i$. In particular, if $\Pr[B_x] \le 2^{-(n+1)}$ for
  every $x \in \{0,1\}^n$, one seed avoids every $B_x$.
\end{lemma}
\begin{proof}
  \leanok
  \uses{lem:encodings-union-count}
  The probabilistic method: the union bound in counting form.
\end{proof}

\begin{lemma}[Acceptance probability is event probability]
  \label{lem:encodings-acceptprob}
  \lean{Complexity.NTM.acceptProb_eq_eventProb, Complexity.NTM.acceptProb_eq_of_le_of_allPathsHaltIn}
  \leanok
  \uses{def:models-accept-prob, def:encodings-eventprob, def:ntm-decides}
  The acceptance probability of $N$ on $x$ with clock $T$ is the event
  probability of the set of accepting choice sequences of length $T$. If, on
  every input $x$, every trace of $N$ halts within $T(|x|)$ steps, and
  $T \le T'$ pointwise, then for every $x$ the acceptance probabilities with
  clocks $T'(|x|)$ and $T(|x|)$ agree.
\end{lemma}
\begin{proof}
  \leanok
  The first statement is definitional. For the second, halted traces are
  fixed and every short choice sequence has the same number of extensions.
\end{proof}

\begin{lemma}[Bridge to Mathlib densities]
  \label{lem:encodings-eventprob-dens}
  \uses{def:encodings-eventprob}
  $\Pr[E]$ equals Mathlib's density \texttt{Finset.dens} of $E$ in
  $\{0,1\}^T$, and the uniform average in Lemma~\ref{lem:encodings-markov} is
  Mathlib's \texttt{Finset.expect}, so the laws above follow from Mathlib's
  general finite-probability API.
\end{lemma}

\section{Repetition of probabilistic machines}

\begin{definition}[Fixed-time repetition]
  \label{def:encodings-repeat}
  \lean{Complexity.NTM.repeatAtTime, Complexity.NTM.repeatAtTimeSteps, Complexity.NTM.repeatAcceptEvent, Complexity.repeatRandomSeed}
  \leanok
  \uses{def:ntm, def:encodings-majority}
  For a machine $N$ with $n$ work tapes, $\mathrm{rep}(N, k, T)$ has $k(n+1)$
  work tapes and performs $k$ sequential runs of $N$ for exactly $T$ steps each,
  each run on a fresh bank of $n + 1$ work tapes (the last simulating the output
  tape), rewinding the input tape between runs, and finally writes the strict
  majority of the $k$ verdicts. Every run occupies exactly $2T + 2$ transitions
  whether or not $N$ halts early, so the total time is $2 + k(2T + 2)$ and the
  choice bits used by the simulated runs sit at fixed positions. The compact
  seed of a choice sequence is the $kT$ bits at those positions, and the
  accepting event of $N$ on $x$ is the set of accepting choice sequences in
  $\{0,1\}^T$.
\end{definition}

\begin{theorem}[Repetition is correct pathwise]
  \label{thm:encodings-repeat-correct}
  \lean{Complexity.NTM.repeatAtTime_trace_correct}
  \leanok
  \uses{def:encodings-repeat}
  Suppose every trace of $N$ on $x$ has halted after $T$ steps. Then for every
  choice sequence $r$ of length $2 + k(2T + 2)$, the trace of
  $\mathrm{rep}(N, k, T)$ on $x$ ends halted with output head at cell $1$ and
  output cell $1$ equal to $\mathrm{maj}_E$ of the compact seed of $r$, where
  $E$ is the accepting event of $N$ on $x$.
\end{theorem}
\begin{proof}
  \leanok
  Invariants for the setup, simulation, rewind, and finish phases of each
  stride.
\end{proof}

\begin{theorem}[Exact acceptance probability of repetition]
  \label{thm:encodings-repeat-prob}
  \lean{Complexity.NTM.repeatAtTime_acceptProb_eq_eventProb}
  \leanok
  \uses{def:encodings-repeat, def:models-accept-prob, def:encodings-eventprob}
  Under the same hypothesis, the acceptance probability of
  $\mathrm{rep}(N, k, T)$ on $x$ with clock $2 + k(2T + 2)$ equals
  $\Pr_{s \in \{0,1\}^{kT}}[\mathrm{maj}_E(s) = 1]$.
\end{theorem}
\begin{proof}
  \leanok
  \uses{thm:encodings-repeat-correct, lem:encodings-acceptprob}
  Acceptance factors through the compact seed, and the administrative choice
  bits cancel because every compact seed has the same number of extensions.
\end{proof}

\begin{theorem}[Amplification by repetition]
  \label{thm:encodings-repeat-amplification}
  \lean{Complexity.NTM.repeatAtTime_acceptProb_ge_one_sub_two_pow, Complexity.NTM.repeatAtTime_acceptProb_le_two_pow}
  \leanok
  \uses{def:encodings-repeat, def:models-accept-prob}
  Suppose every trace of $N$ on $x$ has halted after $T$ steps, and let
  $K = 12s + 1$. If $N$ accepts $x$ with probability at least $2/3$ (clock
  $T$), then $\mathrm{rep}(N, K, T)$ accepts $x$ with probability at least
  $1 - 2^{-s}$ (clock $2 + K(2T + 2)$). If $N$ accepts $x$ with probability at
  most $1/3$, then $\mathrm{rep}(N, K, T)$ accepts $x$ with probability at most
  $2^{-s}$.
\end{theorem}
\begin{proof}
  \leanok
  \uses{thm:encodings-repeat-prob, thm:encodings-amplification}
  Combine the exact identity with majority amplification.
\end{proof}

\section{Zero-error randomized computation}

\begin{definition}[Las Vegas algorithms]
  \label{def:encodings-zpp-lv}
  \uses{def:ntm, def:models-accept-prob}
  $\cc{ZPP}_{\mathrm{LV}}$ is the class of languages $L$ with a probabilistic
  machine $N$ and a polynomial $T$ such that every trace on every input $x$
  halts within $T(|x|)$ steps with output either the correct verdict
  $[x \in L]$ or a designated ``don't know'' answer, and the probability of
  ``don't know'' is at most $1/2$. (The textbook expected-polynomial-time
  formulation needs an expected-time machine semantics, which the library does
  not have.)
\end{definition}

\begin{theorem}[Zero error equals one-sided error on both sides]
  \label{thm:encodings-zpp-lv}
  \uses{def:encodings-zpp-lv, def:ZPP, def:RP}
  \notready
  $\cc{ZPP}_{\mathrm{LV}} = \cc{ZPP}$, where the library defines
  $\cc{ZPP} = \cc{RP} \cap \cc{coRP}$.
\end{theorem}
\begin{proof}
  \uses{thm:encodings-repeat-amplification}
  A Las Vegas machine gives an $\cc{RP}$ machine by answering $0$ on ``don't
  know'', and symmetrically a $\cc{coRP}$ machine. Conversely, run the
  $\cc{RP}$ and $\cc{coRP}$ machines and answer only when one of them
  certifies the verdict.
\end{proof}
